\chapter{Frozen Shoulder}

\section*{Definition and Overview}
Frozen shoulder, medically called adhesive capsulitis, is a condition in which the shoulder becomes progressively stiff, painful, and difficult to move, as the joint's connective tissue capsule becomes inflamed and thickened and then contracts. It typically develops in three overlapping phases: a painful freezing phase, a stiff frozen phase in which pain eases but movement remains limited, and a thawing phase of gradual recovery. The cause is not fully understood, but it is more common in people aged 40--60, in women, and in people with diabetes, thyroid disease, or after an arm injury or surgery. Frozen shoulder usually improves on its own over 1--3 years, and treatment focuses on pain relief and regaining movement. This chapter explains the condition, its phases, and its management.

\section*{Epidemiology}
Frozen shoulder affects around 2--5% of the general population, most often people aged 40--60, and is slightly more common in women. The risk is markedly higher in diabetes, affecting up to 20% of people with diabetes at some point, and it is more common in thyroid disease, heart disease, and Parkinson's disease. Around one in five people develops it in the other shoulder later. It can follow an arm injury, surgery, or a period of immobilisation. The condition is self-limiting in most people, but recovery is slow, and a minority is left with some permanent stiffness.

\section*{Aetiology and Pathophysiology}
\begin{itemize}
    \item \textbf{The joint capsule:} the shoulder is surrounded by a capsule of connective tissue; in frozen shoulder this capsule becomes inflamed, thickened, and contracted, restricting movement
    \item \textbf{Inflammation:} an inflammatory process, possibly triggered by injury or by a systemic condition, affects the capsule
    \item \textbf{Contracture:} over time, the inflamed capsule shrinks and adheres to the joint, freezing it in position
    \item \textbf{Diabetes:} the strong link suggests metabolic factors; people with diabetes get more severe, more bilateral, and more treatment-resistant disease
    \item \textbf{Hormonal and autoimmune links:} thyroid disease and some autoimmune conditions increase risk
    \item \textbf{Immobilisation:} keeping the shoulder still after injury or surgery can precipitate it
    \item \textbf{The three phases:} freezing (pain, 2--9 months), frozen (stiffness, 4--12 months), and thawing (recovery, 6--24 months)
    \item \textbf{Why it recovers:} the inflammation and contracture gradually resolve, restoring movement over time
\end{itemize}

\section*{Clinical Features}
\begin{itemize}
    \item Gradual onset of shoulder pain, often worse at night and when lying on the affected side
    \item Progressive stiffness, with difficulty raising the arm, reaching behind the back, and dressing
    \item Movement limited in all directions, both actively and passively (unlike a tendon problem)
    \item The freezing phase: pain dominates and limits movement
    \item The frozen phase: pain eases but the shoulder is stiff and the range is fixed
    \item The thawing phase: movement gradually returns over months
    \item Pain that interferes with sleep and daily activities
    \item The other shoulder can become affected later
\end{itemize}

\section*{Differential Diagnosis}
Other causes of shoulder pain and stiffness must be considered.
\begin{itemize}
    \item Rotator cuff tears and tendinopathy, which cause pain and weakness but preserve passive movement
    \item Shoulder arthritis, with crepitus and X-ray changes
    \item Subacromial bursitis and impingement
    \item Referred pain from the neck (cervical spine) and from the heart in some cases
    \item Calcific tendinitis
    \item The key distinction: in frozen shoulder, passive movement is restricted, unlike most tendon conditions
\end{itemize}

\section*{When to Seek Medical Care}
Shoulder pain or stiffness that persists, worsens, or limits daily activities should be assessed by a doctor. Early diagnosis allows treatment to begin, and other conditions are excluded.

\textbf{Call 000 (or local emergency services) for:}
\begin{itemize}
    \item Sudden severe shoulder pain with chest pain or breathlessness
    \item Pain after a fall with inability to move the arm, which may indicate a fracture
    \item A painful, swollen, or deformed shoulder after injury
    \item Fever with shoulder pain and swelling, which may indicate infection
\end{itemize}

\section*{Diagnosis and Investigations}
\begin{itemize}
    \item \textbf{History and examination:} the pattern of pain and the restriction of both active and passive movement are usually diagnostic
    \item \textbf{X-rays:} exclude arthritis and calcification
    \item \textbf{Ultrasound or MRI:} assess the rotator cuff and, in some cases, show the thickened capsule
    \item \textbf{Blood tests:} for diabetes and thyroid disease, which are associated and can be newly diagnosed
    \item \textbf{Referral:} to a shoulder specialist or physiotherapist for confirmation and management
\end{itemize}

\section*{Management}
\begin{itemize}
    \item \textbf{Pain relief:} simple analgesics and anti-inflammatory medicines, especially for night pain
    \item \textbf{Physiotherapy:} gentle stretching and range-of-motion exercises, tailored to the phase, to maintain and regain movement
    \item \textbf{Corticosteroid injections:} into the shoulder joint early in the condition to reduce pain and speed recovery of movement
    \item \textbf{Heat and cold:} heat before exercise and ice for pain
    \item \textbf{Activity modification:} using the arm within comfort and avoiding aggravating positions
    \item \textbf{Hydrodilatation:} injecting fluid to stretch the capsule, which helps some people
    \item \textbf{Manipulation under anaesthesia:} the surgeon stretches the shoulder while the person is asleep, used when other treatment fails
    \item \textbf{Arthroscopic release:} surgery to divide the contracted capsule, for persistent stiffness
    \item \textbf{Manage associated conditions:} diabetes and thyroid control
\end{itemize}

\section*{Complications}
\begin{itemize}
    \item Prolonged pain and disability affecting work, sleep, and daily life
    \item Loss of work time, especially in people whose jobs need overhead movement
    \item Permanent stiffness in a minority
    \item Complications of injections and surgery, which are uncommon
    \item Depression and reduced quality of life with long-lasting symptoms
\end{itemize}

\section*{Prognosis and Outcomes}
Frozen shoulder is self-limiting in the majority of people, with gradual recovery over 1--3 years, although a minority is left with some permanent restriction, and a significant number experiences at least mild persistent symptoms. Treatment shortens and eases the course: physiotherapy and corticosteroid injections improve pain and movement, and surgery helps those who do not respond. People with diabetes tend to have a longer, more treatment-resistant course. With patience and a structured program, most people regain good use of the shoulder.

\section*{Prevention and Self-Care}
\begin{itemize}
    \item Keep the shoulder moving gently from the start of any shoulder problem
    \item Perform daily range-of-motion exercises, such as pendulum swings and wall walking
    \item Take analgesia for night pain so sleep is protected
    \item Apply heat before stretching and ice afterwards
    \item Avoid prolonged immobilisation of the arm after injury or surgery
    \item Manage diabetes and thyroid disease well
    \item Be patient: recovery takes months, and improvement is gradual
    \item Seek physiotherapy early for a structured program
\end{itemize}

\section*{Sources and Further Reading}
The following reputable sources provide further information for healthcare students and patients seeking reliable, current advice about frozen shoulder.
\begin{itemize}
    \item Healthdirect Australia. \textit{Frozen shoulder.} https://www.healthdirect.gov.au/frozen-shoulder
    \item Sports Medicine Australia. \textit{Frozen shoulder.} https://sma.org.au/
    \item Australian Physiotherapy Association. \textit{Frozen shoulder.} https://australian.physio/
    \item NHS. \textit{Frozen shoulder.} https://www.nhs.uk/conditions/frozen-shoulder/
    \item Mayo Clinic. \textit{Frozen shoulder.} https://www.mayoclinic.org/
    \item MedlinePlus. \textit{Frozen shoulder.} https://medlineplus.gov/
\end{itemize}
